\documentclass[book.tex]{subfiles}

\begin{document}

\chapter{Introduction to Classical Thermodynamics}

\label{chap:classical_thermo}
\noindent\rule{11cm}{0.4pt} \\
\textbf{Prerequisites:} \autoref{chap:simple_physics}, \autoref{chap:newtons_laws}  \\
\textbf{Difficulty Level:} ** \\
\noindent\rule{11cm}{0.4pt}
\vspace{5mm}

In macroscopic systems, particles constantly change between any of the available energy states predicted by classical or quantum mechanics. Therefore, \textit{ignorance is a law of nature} for many particle systems, and this ignorance requires a statistical description of observations and the admittance of ever-present fluctuations. The presence of such microscopic fluctuations, however, does not prevent a system from reaching macroscopic equilibrium. Thus, before we venture into the subject of characterizing fluctuations, it is necessary to understand the concept of \textit{equilibrium}, which is the subject of classical thermodynamics. 

In the first part of this chapter we will review the concept of thermodynamic equilibrium under the view of classical thermodynamics. Using the first and second law of thermodynamics we will derive the relations that defines a state of macroscopic equilibrium. %We will also introduce the concept of ensembles, which is a fundamental tool to characterize macroscopic systems, and how by using lagrangian transformations we can move from one ensemble to another. We finalize using lagrangian transformations to derive the Euler equation.

\section{Heat, Work, and Entropy}

We will denote heat by $Q$, work by $W$ and entropy with $S$. As we will see in the remainder of this part, understanding the relationships between these quantities will be crucial to deepen our understanding of thermodynamics.

\section{First Law of Thermodynamics}

\begin{marginfigure}%
% \includegraphics[width=\linewidth]{figures/Physics/statistical_mechanics/classical_thermodynamics/thermo_2_systems.png}
\includegraphics[width=1.03\linewidth]{figures/Physics/statistical_mechanics/classical_thermodynamics/Internal energies.png}
\caption{The total internal energy of a system containing two subsystems, is the sum of the internal energies of each subsystem}
\label{fig:Fig1}
\end{marginfigure}

The first law of thermodynamics, also known as the conservation of energy principle, provides the basis for studying the relationships among the various forms of energy and energy interactions. The law states that \textit{although energy assumes many forms, the total quantity of energy is constant, and when energy disappears in one form it appears simultaneously in other forms}.

The energy of a system is measured by its internal energy $U$, and is defined as the net sum of the kinetic and potential energies of all the molecules in a macroscopic material. The internal energy is an extensive property, meaning that the internal energy depends linearly on the size of the system (Fig.\ref{fig:Fig1}).


According to the First Law of Thermodynamics any change in the internal energy occurs due to the addition of heat $Q$ and work $W$, according to the relation:

%\textcolor{red}{TODO: What is the type of $Q$? What is the type of $W$?}

\begin{equation}\label{First Law}
dU = \delta Q + \delta W
\end{equation}

The difference in the notation of the change of internal energy $(dU)$ from the addition of heat $(\delta Q)$ or work $(\delta W)$, lies in the fact that internal energy is a state variable, while heat and work are path dependent. However, once heat and work are applied, they are indistinguishable on a macroscale.

%\textcolor{red}{TODO: How can you represent a path-dependent transformation as a type?}

\section{Second Law of Thermodynamics}

The origins of the second law lies in a statement that it is impossible for any self-acting process or machine to produce net flow of heat from a region of low temperature. The law states that an isolated system always proceeds from an ordered to a more disordered state, or else it undergoes no change at all. 

The level of order from a system is measured by the entropy $S$, and according to the second law \textit{the entropy of the universe tends to a maximum} (Rudolf Claussius, 1865). This is stated mathematically as:

\begin{equation}\label{Second Law}
{\left( dS \right)}_{U,V,N} \ge 0
\end{equation}

\noindent where the subscripts $U$, $V$ and $N$ indicates constant internal energy, volume and number of particles, respectively.

\section{Thermodynamics Driving Forces} 

A thermodynamic system can be defined according to their \textit{extensive} variables, \textit{i.e.}, variables that vary linearly with the system size. For example, we define the internal energy $U$ to be a function of extensive variables $S$, $V$, and $N = (N_1,N_2,...,N_r)$, where $r$ is the number of chemical species in the system. Therefore, we have $U = U(S,V,N)$. 
Differential changes in the extensive variables leads to a differential change in the internal energy, such that:
\begin{equation}\label{Internal Energy Change}
dU ={ \left(\frac{\partial U}{\partial S}\right) }_{V,N_1,...,N_r} \hspace{-1cm} dS 
\hspace{0.6cm} + 
{ \left(\frac{\partial U}{\partial V}\right) }_{S,N_1,...,N_r} \hspace{-1cm} dV
\hspace{0.6cm} + 
{ \left(\frac{\partial U}{\partial N_i}\right) }_{S,V,N_{j\neq i}} \hspace{-1cm} dN_i 
\end{equation}

The partial derivatives are \textit{intensive} variables, \textit{i.e.} variables that do not depended on the size of the system. We define these as:

\begin{equation}\label{intensive variables}
{ \left(\frac{\partial U}{\partial S}\right) }_{V,N_1,...,N_r} \hspace{-1cm} \equiv T \hspace{0.5cm} 
{ \left(\frac{\partial U}{\partial V}\right) }_{S,N_1,...,N_r} \hspace{-1cm} \equiv -p\ 
\hspace{0.5cm} 
{ \left(\frac{\partial U}{\partial N_i}\right) }_{S,V,N_{j\neq i}} \hspace{-1cm} \equiv \mu_i
\end{equation}

Where $T$ is temperature, $p$ is pressure, and $\mu_i$ is chemical potential of species $i$. The change in internal energy can then be written as:

\begin{equation}\label{Eq5}
dU = TdS - pdV + \sum_{i=1}^{r} \mu_idN_i
\end{equation}

Since $U$ is a state variable, we can choose any path to find a differential change. Equation (\ref{Eq5}) states that any change in the internal energy of the system, is the result of a quasi-static heat flux ($TdS$), mechanical work ($- pdV$) and chemical work ($\mu_idN_i$) to the system.

Equivalently, we can define the system according to its entropy $S=S(U,V,N)$, resulting in a differential change:

\begin{equation}\label{Internal Energy}
dS = \frac{1}{T}dU + \frac{p}{T}dV - \sum_{i=1}^r\frac{\mu_i}{T}dN_i
\end{equation}

This formulation is useful to discuss the second law of thermodynamics. 

\begin{marginfigure}%
% \includegraphics[width=\linewidth]{figures/Physics/statistical_mechanics/classical_thermodynamics/thermo_closed_isolated.png}
\includegraphics[width=\linewidth]{figures/Physics/statistical_mechanics/classical_thermodynamics/Isolated systems.png}
\caption{Closed, Isolated System containing two subsystem}
\label{fig:Fig2}
\end{marginfigure}

Consider a closed, isolated system which cannot exchange heat, work, or particles with its surroundings (Fig.~\ref{fig:Fig2}). We initially prepare the system such that part of the system (subsystem 1) has parameter values $U^{(1)}$, $V^{(1)}$, and $N^{(1)}$. We similarly define $U^{(2)}$, $V^{(2)}$, and $N^{(2)}$ for the rest of the system (subsystem 2). Making the statement that the entropy is additive, \textit{i.e.} $S^{(1)} + S^{(2)} = S$), we can write:

\begin{align}\label{Entropy Change1}
dS = \frac{1}{T^{(1)}}dU^{(1)} + \frac{p^{(1)}}{T^{(1)}}dV^{(1)} - \sum_{i=1}^r\frac{\mu_i^{(1)}}{T^{(1)}}dN_i^{(1)} + \notag\\ \frac{1}{T^{(2)}}dU^{(2)}+ \frac{p^{(2)}}{T^{(2)}}dV^{(2)} - \sum_{i=1}^r\frac{\mu_i^{(2)}}{T^{(2)}}dN_i^{(2)}  \hspace{0.3cm} 
\end{align}

Since the total system is closed and isolated, we can write $dU^{(2)}=-dU^{(1)}$, $dV^{(2)}=-dV^{(1)}$, and ${dN}_i^{(2)}=-{dN}_i^{(1)}$.

Therefore, we can write:

\begin{align}\label{Entropy Change2}
dS = \left(\frac{1}{T^{(1)}}-\frac{1}{T^{(2)}}\right)dU^{(1)} + \left(\frac{p^{(1)}}{T^{(1)}}-\frac{p^{(2)}}{T^{(2)}}\right)dV^{(1)} \notag\\ - \sum_{i=1}^r\left(\frac{\mu_i^{(1)}}{T^{(1)}}-\frac{\mu_i^{(2)}}{T^{(2)}}\right)dN_i^{(1)} \hspace{2.8cm} 
\end{align}

%\begin{marginfigure}
%\includegraphics[width=\linewidth]{figures/thermo_polystyrene.png}
%\caption{Phase behavior of polystyrene-poly (2-vinylpyridine) diblock co-polymer. Transmission electron micrographs show lamellar phase (fraction polystyrene $f=0.4$), gyroid phase ($f=0.38$), and cylindrical phase ($f=0.35$). The experimental phase diagram is shown in the lower right panel. TODO: Do we need to replace this?}
%\label{fig:Fig3}
%\end{marginfigure}

According to the second law of thermodynamics, any spontaneous process requires $dS \geq 0$. At equilibrium, the entropy will achieve its maximum and hence satisfy the condition $dS=0$ for arbitrary small changes $dU^{(1)}$, $dV^{(1)}$, and $dN^{(1)}$:

\begin{itemize}
\item Thermal equilibrium occurs when $T^{(1)}=T^{(2)}$
\item Mechanical equilibrium occurs when $p^{(1)}=p^{(2)}$
\item Chemical equilibrium (no reactions) occurs when $\mu_i^{(1)}=\mu_i^{(2)}$
\end{itemize}

Notice that although thermal and mechanical equilibrium implies spatial homogeneity of $T$ and $p$, the condition of spatial homogeneity $\mu_i$ should not be interpreted as homogeneous concentration or density of species $i$. %(vapor-liquid equilibrium in a box or a structured fluid like a block copolymer - Fig.~\ref{fig:Fig3}\cite{Schulz1996}).

%Consider the case $p^{(1)}=p^{(2)}$, $\mu_i^{(1)}=\mu_i^{(2)}$, and the partition between subsystems is rigid and impermeable ($dV^{(1)}={dN}_i^{(1)}=0$).

%\begin{itemize}
%\item If $T^{(1)}>T^{(2)}$ the condition $dS\geq 0$ is met only if $dU^{(1)}<0$
%\item If $T^{(1)}<T^{(2)}$ the condition $dS\geq 0$ is met only if $dU^{(1)}>0$
%\item Heat spontaneously flows from high temperature to low. 
%\end{itemize}

%In the case $T^{(1)}=T^{(2)}$, $\mu_i^{(1)}=\mu_i^{(2)}$, and the partition between subsystems is movable but impermeable (${dN}_i^{(1)}=0$).

%\begin{itemize}
%\item If $p^{(1)}>p^{(2)}$ the condition $dS\geq 0$ is met only if $dV^{(1)}>0$
%\item If $p^{(1)}<p^{(2)}$ the condition $dS\geq 0$ is met only if $dV^{(1)}<0$
%\item Inhomogeneities in pressure leads to the spontaneous expansion of high pressure regions into low pressure regions.
%\end{itemize}

In the case $T^{(1)}=T^{(2)}$, $p^{(1)}=p^{(2)}$:

\begin{itemize}
\item If $\mu_i^{(1)}>\mu_i^{(2)}$ the condition $dS\geq 0$ is met only if $dN_i^{(1)}<0$
\item If $\mu_i^{(1)}<\mu_i^{(2)}$ the condition $dS\geq 0$ is met only if $dN_i^{(1)}>0$
\item Matter flows from high chemical potential to low (not always high concentration to low).
\end{itemize}

\subsection{Helmholtz Free Energy}
%Transforming from $S$ to $T$ for the closed, isothermal ensemble gives the Helmholtz free energy:
The Helmholtz free energy is a quantity defined by the following equation

\begin{align}
F &= U-\left(\frac{\partial U}{\partial S}\right)_{V,N} S \\
&= U -TS
\end{align}

which has a differential change:
\begin{align}
dF &= dU -TdS-SdT \\
&=-SdT-pdV + \sum_{i=1}^r\mu_idN_i
\end{align}

\begin{lstlisting}
F: System $\to$ $\mathbb{R}$
\end{lstlisting}

\subsection{Gibbs Free Energy}
The Gibbs free energy is defined by the following quantity
%Transform from $S$ to $T$ and $V$ to $p$ for the closed, isothermal, isobaric ensemble give the Gibbs free energy:


\begin{align}
G &= U - \left(\frac{\partial U}{\partial S}\right)_{V,N}S - \left(\frac{\partial U}{\partial V}\right)_{S,N}V \\
&= U - TS + pV
\end{align}
which has a differential change:
\begin{align}
dG &= dU -TdS - SdT + pdV +Vdp \\
&= -SdT +Vdp + \sum_{i=1}^r\mu_idN_i 
\end{align}

\begin{lstlisting}
G: System $\to$ $\mathbb{R}$
\end{lstlisting}

\section{Euler Equation}
The internal energy is a homogeneous, first-order function, which
requires that the internal energy must increase linearly with the size
of the system. Mathematically, this is written as:
\begin{equation}
U(\lambda S,\lambda V, \lambda N) = \lambda U(S,V,N)
\end{equation}

where $\lambda$ is an arbitrary scaling factor for the system size. If we take the derivative of this equation with respect to $\lambda$, we arrive at:

\begin{equation}
\frac{\partial U(\lambda S,\lambda V, \lambda N)}{\partial \lambda} = U(S,V,N)
\end{equation}

The left-hand side is written as:
\begin{align}
\left(\frac{\partial U(\lambda S,\lambda V, \lambda N)}{\partial (\lambda S)}\right)_{\lambda V, \lambda N} \frac{\partial (\lambda S)}{\partial \lambda} + \left(\frac{\partial U(\lambda S,\lambda V, \lambda N)}{\partial (\lambda V)}\right)_{\lambda S, \lambda N} \frac{\partial (\lambda V)}{\partial \lambda} + \notag \\
\sum_{i=1}^r\left(\frac{\partial U(\lambda S,\lambda V, \lambda N)}{\partial (\lambda N_i)}\right)_{\lambda S, \lambda V, \lambda N_{j \neq i}} \frac{\partial (\lambda N_i)}{\partial \lambda}  = TS - pV + \sum_{i=1}^r\mu_iN_i 
\end{align}

Therefore, the total Legendre Transform equals zero, which is the Euler equation:
\begin{equation}\label{Euler equation}
U = TS -pV +\sum_{i=1}^r\mu_iN_i
\end{equation}


\end{document}